%%%%%%%%%%%%%%%%%%%%%%%%%%%%%%%%%%%%%%%%%%%%%%%%%%%%%%%%%%%%%%%%%%%%%%%
% CHAPTER 1
%
% This chapter is laid out as a complete empirical study, so it doubles
% as a skeleton for a paper-style thesis chapter. Delete the sections
% you do not need; the headings are here to show the architecture, not
% to prescribe it.
%%%%%%%%%%%%%%%%%%%%%%%%%%%%%%%%%%%%%%%%%%%%%%%%%%%%%%%%%%%%%%%%%%%%%%%
\documentclass[../main.tex]{subfiles}

\begin{document}
\newpage
\onehalfspacing
\RaggedRight
\setcounter{chapter}{0}
\thesischapterstyle

\chapter{Chapter One Title}\label{chap1:short-label}

%======================================================================
\section{Introduction}\label{sec:ch1-introduction}

\subsection{Background}
% Open on the phenomenon, with a quantified magnitude and a citation
% establishing that it is first-order. An event or a date alone is not
% a hook.
\lipsum[1]

% Citations: \citep{key} renders as (Author 2020); \citet{key} renders as
% Author (2020). Both come from references.bib.
Prior work documents this pattern \citep{placeholder-article}, and
\citet{placeholder-book} provides the standard treatment.

\subsection{Research Motivation}
% Why the question is open: what the existing evidence cannot settle.
\lipsum[2]

\subsection{Research Questions}
% State them as questions, numbered, each answerable by one exhibit.
\begin{enumerate}
    \item Research question 1.
    \item Research question 2.
    \item Research question 3.
\end{enumerate}

\subsection{Main Findings}
% One paragraph per headline result, each pointing at the table that
% carries it.
\lipsum[3]

\subsection{Research Contributions}
% Name the closest rival studies by author-year with live citations and
% state the delta. Naming a literature is not naming a rival.
\lipsum[4]

\subsection{Chapter Structure}
The remainder of this chapter is organised as follows.
Section~\ref{sec:ch1-literature} reviews the related literature.
Section~\ref{sec:ch1-theory} develops the theoretical framework.
Section~\ref{sec:ch1-hypotheses} sets out the hypotheses.
Section~\ref{sec:ch1-data} describes the data and empirical model.
Section~\ref{sec:ch1-results} reports the results.
Section~\ref{sec:ch1-conclusion} concludes.

%======================================================================
\section{Literature Review}\label{sec:ch1-literature}

% Density follows function: the closest rivals get paragraphs,
% background gets a clause. Every cited work should do one of four
% jobs -- establish a premise, define the gap, supply a method
% precedent, or be the rival the chapter must beat.

\subsection{First Strand of Literature}
\lipsum[5]

\subsection{Second Strand of Literature}
\lipsum[6]

\subsection{Synthesis: Research Gaps}
\lipsum[7]

%======================================================================
\section{Theoretical Framework}\label{sec:ch1-theory}

\subsection{Core Framework}
\lipsum[8]

\subsection{Mechanism Specification}
\lipsum[9]

\subsection{Conceptual Model}

Figure~\ref{fig:ch1-framework} summarises the conceptual model.

\begin{figure}[p]
\centering
\begin{tikzpicture}[
    >=stealth,
    node distance = 2.2cm,
    box/.style = {rectangle, draw=black, thick, rounded corners,
                  minimum width=3cm, minimum height=1cm,
                  align=center, font=\small}
]
\node[box, fill=gray!12] (x)   {Variable of interest\\$X$};
\node[box, fill=gray!12, right=of x] (m) {Mechanism\\$M$};
\node[box, fill=gray!12, right=of m] (y) {Outcome\\$Y$};
\node[box, below=1.4cm of m] (z) {Moderator\\$Z$};

\draw[->, thick] (x) -- node[above, font=\scriptsize] {H1} (m);
\draw[->, thick] (m) -- node[above, font=\scriptsize] {H2} (y);
\draw[->, thick, dashed] (z) -- node[right, font=\scriptsize] {H3} (m);
\end{tikzpicture}
\caption{Conceptual framework.}
\label{fig:ch1-framework}
\end{figure}

%======================================================================
\section{Hypothesis Development}\label{sec:ch1-hypotheses}

% Each hypothesis should trace a chain: established finding (cited)
% -> the step this chapter adds -> the prediction. State the expected
% sign for each outcome separately, not in general.

\subsection{Main Hypothesis}
\lipsum[10]

\begin{quote}
\textbf{Hypothesis 1.} \textit{State the predicted relationship and its
sign here.}
\end{quote}

\subsection{Mechanism Hypotheses}
\lipsum[11]

\begin{quote}
\textbf{Hypothesis 2.} \textit{State the predicted mechanism here.}
\end{quote}

%======================================================================
\section{Data and Model}\label{sec:ch1-data}

\subsection{Data and Sample}
% Name every source, the coverage period, and the unit of observation.
\lipsum[12]

\subsection{Variable Construction}

\subsubsection{Dependent Variable}
% Every filter and threshold should name its precedent, or be labelled
% explicitly as this chapter's own choice.
\lipsum[13]

\subsubsection{Variable of Interest}
\lipsum[14]

\subsubsection{Control Variables}
\lipsum[15]

\subsection{Research Model}

The baseline specification is estimated as

\begin{equation}\label{eq:ch1-baseline}
    Y_{i,t} = \beta_0 + \beta_1 X_{i,t}
              + \sum_{k} \gamma_k \, \text{Control}^{k}_{i,t}
              + \mu_i + \lambda_t + \varepsilon_{i,t},
\end{equation}

\noindent where $Y_{i,t}$ is the outcome for unit $i$ in period $t$,
$X_{i,t}$ is the variable of interest, $\mu_i$ and $\lambda_t$ denote
unit and period fixed effects, and $\varepsilon_{i,t}$ is the error
term. Equation~\eqref{eq:ch1-baseline} is estimated by \ac{OLS}, with the
\ac{SE} clustered at the [LEVEL] level. The
coefficient of interest is $\beta_1$; Hypothesis~1 predicts
$\beta_1 > 0$.

%======================================================================
\section{Empirical Results}\label{sec:ch1-results}

\subsection{Descriptive Statistics}

Table~\ref{tab:ch1-descriptive} reports the descriptive statistics.

% [p] sends the table to a float page rather than squeezing it beside the
% text -- the empirical finance convention, and what a full-width
% regression table needs anyway.
%
% Note what LaTeX does with it: a float page holds as many floats as fit,
% so the two demo tables and the framework figure below all land on one
% page. Real regression tables are large enough that each takes a page of
% its own, so this does not happen in a finished thesis. If yours do pack
% together, they are small enough to sit inline -- use [htbp] for those.
%
% For a journal submission the convention is the opposite: markers in the
% text, exhibits at the end. See the float-placement note in main.tex --
% switch it there rather than moving anything by hand.
\begin{table}[p]
    \centering
    \caption{Descriptive statistics}
    \label{tab:ch1-descriptive}
    \begin{tabular}{lccccc}
        \toprule
        Variable & N & Mean & SD & Min & Max \\
        \midrule
        $Y$          & [N] & [.] & [.] & [.] & [.] \\
        $X$          & [N] & [.] & [.] & [.] & [.] \\
        Control 1    & [N] & [.] & [.] & [.] & [.] \\
        Control 2    & [N] & [.] & [.] & [.] & [.] \\
        Control 3    & [N] & [.] & [.] & [.] & [.] \\
        \bottomrule
    \end{tabular}

    \vspace{1mm}
    {\footnotesize
    This table reports summary statistics for the estimation sample over
    [PERIOD]. All continuous variables are winsorised at the
    [X]th and [100-X]th percentiles. Variable definitions are given in
    Appendix~\ref{app:variables}.
    }
\end{table}

\subsection{Baseline Results}

Table~\ref{tab:ch1-baseline} reports estimates of
Equation~\eqref{eq:ch1-baseline}.

\begin{table}[p]
    \centering
    \caption{Baseline results}
    \label{tab:ch1-baseline}
    \begin{tabular}{lcccc}
        \toprule
                     & (1) & (2) & (3) & (4) \\
                     & $Y$ & $Y$ & $Y$ & $Y$ \\
        \midrule
        $X$          & [.]$^{***}$ & [.]$^{***}$ & [.]$^{**}$  & [.]$^{**}$ \\
                     & ([.])       & ([.])       & ([.])       & ([.])      \\
        Control 1    &             & [.]         & [.]         & [.]        \\
                     &             & ([.])       & ([.])       & ([.])      \\
        Control 2    &             & [.]         & [.]         & [.]        \\
                     &             & ([.])       & ([.])       & ([.])      \\
        Constant     & [.]         & [.]         & [.]         & [.]        \\
                     & ([.])       & ([.])       & ([.])       & ([.])      \\
        \midrule
        Controls     & No  & Yes & Yes & Yes \\
        Unit FE      & No  & No  & Yes & Yes \\
        Period FE    & No  & No  & No  & Yes \\
        Observations & [N] & [N] & [N] & [N] \\
        Adj.\ $R^2$  & [.] & [.] & [.] & [.] \\
        \bottomrule
    \end{tabular}

    \vspace{1mm}
    {\footnotesize
    This table reports OLS estimates of Equation~\eqref{eq:ch1-baseline}.
    Standard errors clustered at the [LEVEL] level are in parentheses.\\
    $^{***}$ $p<0.01$, $^{**}$ $p<0.05$, $^{*}$ $p<0.1$.
    }
\end{table}

\subsection{Identification and Endogeneity}
% Instrumental variables, difference-in-differences, propensity score
% matching, or whatever the design requires. Mechanism evidence comes
% after identification, never before.
\lipsum[16]

\subsection{Mechanism Tests}
\lipsum[17]

\subsection{Heterogeneity Analysis}
% Moderators enter as interactions on dichotomised variables in the
% baseline; saturated high-order interactions belong in robustness.
\lipsum[18]

\subsection{Robustness Tests}
\lipsum[19]

%======================================================================
\section{Conclusion}\label{sec:ch1-conclusion}

\subsection{Summary of Findings}
\lipsum[20]

\subsection{Theoretical Implications}
\lipsum[21]

\subsection{Practical Implications}
\lipsum[22]

\subsection{Limitations and Future Research}
\lipsum[23]

\end{document}
